\chapter{Outro}

Having discussed open questions in the outro section of each chapter, I guess all that really remains for this overall outro is to muse on the legacy of this thesis.
I suppose I am reasonably happy with the second half of it -- I think Pauli webs will be a very useful tool for qumodes just as they have been for qubits, and in fact I think CV quantum information is arguably more in need of graphical tools that its finite dimensional counterpart, just because it's a good deal more complicated!
Though maybe that's just my prejudice (i.e.\ fear of calculus!) coming through.
The wider usefulness of this CV stuff probably depends on the success (or not) of CV quantum computing/error correction as a whole, and being a humble qubit fan at heart that's not something I reckon I'm qualified to speculate on.

I have more mixed feelings about the first half of this thesis.
On the one hand, I hope \Cref{ch:dynamic-stabilizer-codes} ends up being a nice resource for learning about dynamic stabilizer codes, filling a bit of a gap in the literature.
Furthermore, the work done in \Cref{ch:floquetifying} was the first time in my PhD I really got the research bug, staying up late drawing silly ZX-diagrams and being excited to share progress the next day.
This was a welcome change from the first few months of my PhD where I wanted to become an LDPC code king, but in hindsight went about it in rather the wrong way and ended up feeling quite lonely and disenchanted.
On the other hand, I think it's already clear that this chapter won't stand the test of time!
Firstly, the fact that one can reinterpret a static stabilizer code as a dynamic stabilizer code turned out not be as interesting as first hoped -- as we've now seen, \textit{any} circuit implementing the static code can be viewed as a dynamic stabilizer code in its own right.
Secondly, we made no attempt to think about behaviour of our new dynamic stabilizer code under noise, and indeed it turned out that our proposed code was laughably bad in a practical sense.

The saving grace seems to be that this work played a small part in Rodatz et al.\ going on to do what I reckon is genuinely excellent work on fault-equivalence~\cite{rodatz2024floquetifying,rodatz2025faulttoleranceconstruction,rüsch2025completenessfaultequivalenceclifford}.
It's very likely they'd have got there anyway given that we didn't do much more than PsiQuantum had already done in \Ccite{Bombin_2024} for the surface code, but if that ends up being the legacy of these four years then I'm more than okay with it.
It's been a lovely four years, in any case.

\pagebreak

\textit{Acknowledgements.} First, I wish to thank the head of the Eisert empire, Jens Eisert himself, for taking me on (despite him being a ZX-calculus skeptic at the time), providing us with total academic freedom, shielding us from a lot of the boring stuff, and being very generous re: attending workshops, summer schools, conferences, and so on.
Like any good zealot, I am happy to have played a small part in changing his mind on the ZX-calculus over the last few years.
I'm very grateful to our admin team Claudia Thomas and Felix Witte too, for keeping our frankly monstrous group on an even keel throughout this time.

Big thanks to my QEC office mates PJ Derks, Ansgar Burchards and Johannes Frank, for making the day-to-day fun.
I'm grateful to everyone I've collaborated with over the course of my PhD -- Markus Kesselring, Julio Magdalena de la Fuente, Oscar Higgott, Ben Brown, Josias Old, Razin Shaikh, Ben Cichos and Henri Wegener.
More widely, I'm grateful to everyone I've had discussions with over the years -- too many to name, but I'd like to give a particular shoutout to everyone in the QEC subgroup here in Berlin, to Dan Litinski and Fernando Pastawski for the webs, and to Ben Rodatz.

By far the most important event during my PhD was the IBM QEC Summer School in 2022, where I learned loads and made a ton of close friendships -- again, too many to name.
For this, I'm hugely thankful to Drew Vandeth and Anamaria Rojas for the inception and organisation, and to all the lecturers who gave their time.
